% Proposed B162, first saved33 source. No central edits.
\subsection*{The full topology of the double and its actual continuous free kernel}
\statement{B.162}{Proved}{The actual unfilled double has full induced topology, while its infinite discrete fold kernel does not; the actual closed kernel has a continuous pointed-coset free model}
\label{stmt:B162}
Use B.161. The groups $D_+$ and $D_0$ have their full induced profinite
topologies in their actual generic ambient groups. Write bars for the
corresponding actual closures. If
\[
 X=\widehat J/\widehat S,\qquad *=\widehat S,
\]
and $U$ is the free profinite group on the \emph{continuous pointed}
space $(X,*)$, then
\begin{equation}\label{eq:B162-model}
 \overline{D_0}=\widehat{D_0}
       =U\rtimes\widehat\lambda(\widehat J),\qquad
 \overline R=U.
\end{equation}
The notation means continuous pointed families of elements, not the
abstract profinite completion of the discrete group in B.161.
In fact $R$ does not have full induced topology in $D_0$, and
$\widehat R\to\overline R$ has nontrivial kernel.
\begin{proof}
\emph{Extend every finite double target by finite permutations.}
In the legitimate pre-node quotient of B.161 set
$B_0=S\times C_2(\tau)\le Q_c$. The subgroup $H$ is a retract of
$Q_c$ via its isomorphic projection onto $S$. Killing $C,f$ gives a
retraction onto $B_0$. Hence both have full induced topology.
B.153 gives
\[
 \widehat H\cap\widehat S=\widehat{\langle m\rangle}.
\]
Since $H$ is even and the even part of $B_0$ is $S$, this implies
$\widehat H\cap\widehat B_0=\widehat{\langle m\rangle}$.

Let $\psi:D_+\to F$ be any finite target. In
$F'=(F\times F)\rtimes C_2(\omega)$, where $\omega$ swaps the
coordinates, put
\[
 z=(\psi(A),1)\omega,\qquad j=(\psi(\tau),\psi(\tau)),
 \qquad a_0=zj.
\]
Then $z^2=\operatorname{diag}\psi(A)$, $j^2=1$, and $a_0$ commutes
with $\operatorname{diag}\psi(m)$. Moreover
$(a_0j)^2=\operatorname{diag}\psi(A)$.
Let $K_B=\ker(\psi|_{B_0})$. The compact set
\[
 \widehat B_0\setminus
       (\widehat{\langle m\rangle}\,\widehat K_B)
\]
is disjoint from $\widehat H$. Thus there is an open normal subgroup
$V\trianglelefteq\widehat Q_c$ such that $V\cap\widehat B_0\le
\widehat K_B$ and the image intersection satisfies
\[
 \beta(\overline H\cap\overline B_0)
       \subseteq\operatorname{diag}\psi(\langle m\rangle),
 \quad \beta:\overline B_0\to F',\quad
 \overline Q_c=\widehat Q_c/V.
\]
Indeed the decreasing closed sets $\widehat HV$ intersect in
$\widehat H$; compact separation gives the first condition, and
the induced subgroup topology supplies the second kernel refinement.

Form the finite induced set
\[
 \Omega=\overline Q_c\times_{\overline B_0}F',\qquad
 (qb,t)\sim(q,\beta(b)t),\qquad \Omega_0=\{[1,t]:t\in F'\}.
\]
The left $\overline Q_c$ action is well-defined. On
$\overline H\Omega_0$ define $P_a([h,t])=[h,a_0t]$, and make it
the identity on the complement. Two representatives differ by an
element of $\overline H\cap\overline B_0$, whose image commutes with
$a_0$. Thus this is a well-defined invertible permutation commuting
with the whole $\overline H$ action. It gives a finite permutation
target of $Q_c*_{H}(H\times\mathbb Z_a)=E_{\rm pre}$.

The actual $D_+$ preserves $\Omega_0$. On that fiber, $S,\tau$ act
by $\operatorname{diag}\psi$, and $A_0=(a\tau)^2$ acts by
$(a_0j)^2=\operatorname{diag}\psi(A)$. Therefore the kernel of this
ambient finite action restricts to a subgroup of $\ker\psi$.
Composition with the original generic deletion proves full topology
of $D_+$. Finite-index core refinement proves the same for its
index-two subgroup $D_0$ in the ordinary ambient group. Evaluation
was used only on the invariant fiber; no map of the whole ambient
group to the original finite codomain $F$ is asserted.

\emph{Construct the continuous pointed group, with both inverse maps.}
For each open normal $V\trianglelefteq\widehat J$, put
\[
 X_V=\widehat J/(\widehat S V),\qquad
 U_V=\widehat{F(X_V\setminus\{*\})},\qquad U=\varprojlim_V U_V.
\]
The transition maps fold generators with equal cosets and send a
generator to $1$ if its image is the base coset. They are surjective.
The compact coset space $X$ is the inverse limit of $X_V$.
Every continuous pointed map from $X$ to a finite group factors
through some $X_V$: its pullback to $\widehat J$ is locally constant
and right $\widehat S$-invariant, and a common open normal subgroup
refines its finitely many clopen fibers. Extension to $U_V$ and then
finite-target inverse limits proves the asserted universal property.
The family of generators topologically generates $U$, because it
does so in every finite-level free profinite group.

The formulas
\[
 \theta_g(X_{h\widehat S})=X_{g\widehat S}^{-1}X_{gh\widehat S}
\]
give compatible continuous actions at every finite level, hence a
jointly continuous action on $U$. The two maps
$j\mapsto j$ and $j\mapsto X_{j\widehat S}j$ into
$U\rtimes\widehat J$ agree on $\widehat S$. Conversely the actual
cocycle of the two copies in $\widehat{D_0}$ is a continuous pointed
family on $X$. The universal property gives the reverse map. The
identities in B.161 show that these maps respect the action and are
inverse on every marked generator. This proves the completed double
model, without identifying a completion of an infinite-index kernel.
The completed fold kernel is $U$; it equals $\overline R$ by right
exactness, or by correcting an approximating net using the fold section.
Full topology of $D_0$ identifies this with the actual ambient closure.
The swapping involution still satisfies $\delta(X_x)=X_x^{-1}$.

\emph{The discrete kernel has strictly finer finite quotients.}
The character of the free group $R$ sending $X_{AS}$ to $1\in C_2$
and every other basis letter to $0$ cannot be refined by any finite
$D_0$ target. Indeed choose a positive integer $n$ divisible by the
orders of both images of $A$ in its two copies. Then
\[
 X_{A^{n+1}S}\longmapsto X_{AS}
\]
in that target. The cosets $S,AS,A^{n+1}S$ are distinct by the
$A$-exponent character of $J$, whereas the chosen character separates
the latter two free letters. Thus even kernel refinement fails.
If $\widehat R\to\overline R$ were injective, it would be a compact
isomorphism, so this character would be continuous on $\overline R$.
An open ambient normal subgroup refines every continuous finite target
of a closed subgroup, contradicting the preceding argument. This also
explains why the continuous pointed model cannot be replaced by
$\widehat R$.
\end{proof}
The finite permutation extension is the complete root/DE/CE argument
audited during the outage and now bound to B.139's restored pre-node input.
No separability or virtual-retraction assertion about the discrete
infinite kernel is included.
